\section{Relation to Olde Loohuis and Venema}\label{app:olv}

Olde Loohuis and Venema \cite{OldeLoohuisVenema2010} and this paper ask the same basic question with different observations. Each takes the value of a multiplayer evaluation game, turns agreement at every valuation into an equational theory, and asks how to decide and axiomatize it. They observe who can force a win; we observe who wins in some subgame-perfect equilibrium. This appendix records how the two settings correspond, which of their questions the paper answers, and which remain open here. Section numbers below refer to their paper.

\subsection{The two settings}

\begin{center}
\begin{tabular}{p{0.17\linewidth}p{0.37\linewidth}p{0.37\linewidth}}
\toprule
& Olde Loohuis and Venema & This paper\\
\midrule
Value & the players who can force a win & the players who win in some pure SPE\\
Leaves & any set of winners, possibly none & exactly one winner\\
Carrier & all subsets of $I$ & the nonempty subsets of $I$\\
Choice by $r$ & join in $r$'s coordinate, meet in the others & $\{r\}$ absorbs; otherwise union\\
Role change & negations $\neg_{ij}$, each swapping two roles & renamings $\pi\in\Sym(I)$\\
Two players & Belnap's four-valued logic & strong Kleene logic\\
\bottomrule
\end{tabular}
\end{center}

Their algebra $O_I$ \cite[Definition~3.20]{OldeLoohuisVenema2010} is a power of the two-element lattice, with one coordinate per player, and each choice acts coordinatewise as join or meet. By \cref{cor:no-boolean-product}, the winner algebra embeds in no such product, so the two algebras differ structurally, not only in what they observe. On games with one winner at each leaf, a value is $\{i\}$ exactly when $i$ can force a win, and an ambiguous value means that nobody can. Identifying all ambiguous values therefore gives the choice operations of $O_I$ restricted to its empty and singleton values (\cref{game:prop:obs} and the discussion following it). With two players, the winner algebra is the subalgebra of theirs in which $\{a,b\}$ is the value where neither player can force a win (\cref{app:two-players}). Their negations swap two roles, and these transpositions generate all renamings.

\subsection{Their questions and ours}

\paragraph{Deciding equations.} Equivalence is coNP-complete in both settings. Their hardness proof lets two owners act as a distributive lattice \cite[Theorem~5.50]{OldeLoohuisVenema2010}, using the hardness of the word problem for distributive lattices \cite{BloniarzHuntRosenkrantz1984}. Owned choice is not distributive, and individual forcing does not define possible winners, so that reduction does not transfer. The paper instead locates where hardness starts (\cref{sec:complexity}) and shows that pairs suffice to refute equations (\cref{sec:pair-tests}).

\paragraph{Axiomatization.} Their equations have a finite basis \cite[Definition~3.24]{OldeLoohuisVenema2010}, and $O_I$ generates their variety. Owned choice has no finite basis, nor one in boundedly many variables (\cref{core:thm:core-nfb,full:thm:nonfinite-basis}). Finite bases exist for fragments: owned choice with union (\cref{core:prop:choice-union-basis}), reset (\cref{reset:thm:basis}), reset with one guarded choice (\cref{owner:thm:basis-main}), and the quotient that identifies all ambiguous values (\cref{core:thm:atomic-basis}).

\paragraph{Representation.} They prove a representation theorem analogous to Stone's \cite[Theorem~3.28]{OldeLoohuisVenema2010} and a dual adjunction with sets. We have no analogue.

\paragraph{One player's view.} Their player $i$ reads a value through one bit, whether $i$ can force a win, and these bits determine the value \cite[Proposition~3.26]{OldeLoohuisVenema2010}, so their algebra is the product of its views. Here a player reads $K+1$ states, since it must recognize every other player's forced win (\cref{view:thm:view}), and the views give only a subdirect product (\cref{view:cor:product}).

\paragraph{Expressiveness.} They show that their language is not functionally complete and leave the relative expressive power of sets of connectives aside. The paper determines the idempotent operations (\cref{gen:thm:core}) and the binary atom-preserving ones (\cref{reset:thm:binary}), defines union from owned choices (\cref{game:thm:union}), and shows that no finite family of relations characterizes the core's term operations (\cref{core:thm:not-finitely-related}).

\paragraph{Single-winner valuations.} They introduce valuations that give each proposition letter one winner, note that a formula can then have no player who forces a win, and leave the logic of these valuations open \cite[Definition~2.12]{OldeLoohuisVenema2010}. One winner at each leaf is our standing assumption, and a game has an ambiguous value exactly when nobody can force a win in it. The basis of \cref{core:thm:atomic-basis} axiomatizes the choice operations of $O_I$ on its empty and singleton values, but its variables may take the empty value, so it is not the logic of single-winner valuations itself. Equation validity stays coNP-complete when every variable takes an atom: the two-owner reduction in the proof of \cref{thm:raw-pairs} refutes the equation of a satisfiable formula with an atomic valuation, and an unsatisfiable formula gives an equation valid at every valuation.

\paragraph{One player's possible win.} They decide whether a player can win under some valuation in linear time, by evaluating at the valuation in which every player wins every letter, and observe that single-winner valuations make the problem NP-complete \cite[Theorem~5.49]{OldeLoohuisVenema2010}. The first part holds here too.

\begin{proposition}[One player's possible win]\label{olv:prop:sat}
For every term $t$ of $\mathbf A^{\rm sym}$ and every role $k$, some valuation puts $k$ in the value of $t$ if and only if the valuation assigning $I$ to every variable does. So this can be decided in linear time at fixed $K$.
\end{proposition}
\begin{proof}
Owned choice is monotone for inclusion: if $S\subseteq S'$, then $S\ow r T\subseteq S'\ow r T$. Indeed, if $S'\ow r T=\{r\}$, then $S'=\{r\}$ or $T=\{r\}$, and in both cases $S\ow r T=\{r\}$. Otherwise $S'\ow r T=S'\cup T$, which contains $S\ow r T$, whether that is $S\cup T$ or $\{r\}$ with $r\in S$. Renamings are monotone and winner leaves are constant, so every term operation is monotone, and its value is largest when every variable is $I$.
\end{proof}

\subsection{Their directions for further research}

\paragraph{Rational play.} Their Section~6 proposes rationality, and common knowledge of rationality, as a refinement, and names backward induction and equilibria as solution concepts to relate to. They sketch a clause for a choice by $j$: if $j$ can force a win in exactly one branch, another player wins the choice exactly when it wins that branch; otherwise the original semantics applies. With one winner at each leaf, this clause changes nothing.

\begin{proposition}[Their rationality clause]\label{olv:prop:rational}
On games with one winner at each leaf, the rationality clause of Olde Loohuis and Venema \cite[Section~6]{OldeLoohuisVenema2010} assigns each player exactly the games in which it can force a win.
\end{proposition}
\begin{proof}
At most one player can force a win in such a game, since two forcing strategies played together would reach a leaf with two winners. Induct on the game. At a choice by $j$, the clause departs from forcing for a player $i\ne j$ only when $j$ can force a win in exactly one branch, and it then asks $i$ to win that branch, which $i$ cannot do. Forcing gives the same answer, since $i$ would have to force a win in both branches. For $j$ itself the clause is the original one.
\end{proof}

So the guaranteed winners under their clause are the players who can force a win, and possible SPE winners are what yield a new algebra.

\paragraph{Uncertainty.} They suggest imperfect and incomplete information. The sealed rounds of \cref{app:sealed} add simultaneous hidden moves; with three players, no compositional semantics with finitely many values tracks even who might win (\cref{se:thm:infinite}).

\paragraph{Coalitions.} They suggest operators indexed by groups of players, or negations that form coalitions, possibly by non-injective maps of roles. Here a coalition that shares one objective behaves as a single role, and a team that wins when any member does has the same identities as the core for three or more roles (\cref{context:prop:team-refinement}); and \cref{context:prop:observable-coalitions} determines which coalition powers the winner operation records. Renamings are always permutations here.

\paragraph{Fixpoints.} They propose a multiplayer $\mu$-calculus, whose plays may be infinite. Our games are well founded (\cref{game:prop:choice-support}); infinite plays are not treated.

\paragraph{Role switches.} They ask whether the role switches are determined by the other operations, as a Boolean algebra is determined by its bounded lattice reduct. For the winner algebra itself they are.

\begin{proposition}[Renamings are determined]\label{olv:prop:renaming}
Let $K\geq3$ and $\pi\in\Sym(I)$. If a bijection $\varphi$ of the values satisfies $\varphi(S\ow r T)=\varphi(S)\ow{\pi(r)}\varphi(T)$ for all $r$, $S$, and $T$, then $\varphi(S)=\pi(S)$ for every value $S$.
\end{proposition}
\begin{proof}
The atom $\{s\}$ is the only value absorbing for $\ow s$: if $Z$ absorbs, then $Z=\{s\}\ow s Z=\{s\}$. Since $\varphi$ is onto, $\varphi(\{r\})$ absorbs for $\ow{\pi(r)}$, so $\varphi(\{r\})=\{\pi(r)\}$. For distinct $r,s,t$, \cref{game:thm:union} gives $\varphi(S\cup T)=\varphi(S)\cup\varphi(T)$, because $\pi(r),\pi(s),\pi(t)$ are again distinct. A value is the union of its atoms.
\end{proof}

\subsection{Open questions}

\begin{itemize}
\item The logic of single-winner valuations: a basis, or a proof that none is finite, for the equations that hold when every variable takes an atom.
\item To do: the complexity of deciding whether a player can win under some atomic valuation, which they find NP-complete in their setting. The reductions of \cref{sec:complexity} use atomic countervaluations, so NP-hardness is expected here as well.
\item A representation or duality theorem for the varieties generated by the winner algebras, as they give for theirs.
\item Coalitions formed by non-injective maps of roles, which change incentives and so are not renamings of values.
\item Infinite plays, and multiplayer fixpoint operators.
\end{itemize}
